本文提出了 KCH-MedRank，一个面向生物医学证据检索的知识约束学习排序框架。核心发现是医学知识约束将超图结构从噪声转化为信号。没有领域标注的原始超图拓扑会损害排序。加入 MeSH、实体和 PrimeKG 约束后，同一架构产生可解释的、统计显著的改进。

在 BioASQ 上，KCH-MedRank 在 Recall@10 上显著优于真实 MedCPT Cross-Encoder 基线（$+0.0157$，$p=0.0076$），重排序快 $18.6$ 倍。扁平生物医学知识特征承担了大部分聚合收益。超图在 $k{=}5$ 上贡献了针对性的早期排序收益，集中在查询与证据之间缺少直接 MeSH 重叠的问题上。BEIR NFCorpus 确认了监督检索特征重排序的稳健性，PubMedQA 诊断显示证据排序在受控条件下影响下游答案选择。

结构消融直接证明了领域概念标注对基于图的生物医学检索是必要的。从词汇或嵌入相似度建边的通用 GraphRAG 和 HyperGraphRAG 流水线面临引入结构噪声而非有用信号的风险。可解释性分析将 648 条救回金标准段落追溯到三种知识机制，MeSH 层级路径占救回的 75.9\%。错误分析揭示了生物医学概念归一化，特别是处理同义词和缩写，是最值得投入的改进方向。

KCH-MedRank 是一个可解释、高效、知识驱动的证据重排序模块。代码和已处理特征文件将在发表后发布。
